%% cap02-shorter-hesus-speaks.tex — Marco Teórico y Contextual

\chapter{Marco Teórico y Contextual}
\label{cap:marco-teorico}

\begin{quote}\small\textit{
Nota al lector: En este capítulo construí el andamiaje conceptual de la investigación. Partí de lo que la inteligencia artificial significa para la educación, revisé los marcos internacionales que han moldeado el discurso global, y después abordé cómo se definen y cómo viajan las políticas educativas entre países. La tradición de la educación comparada proporcionó la estructura de análisis, y las herramientas de procesamiento de lenguaje natural, la capacidad de operar sobre textos en cuatro idiomas. El capítulo cierra con el marco analítico de nueve dimensiones que organiza la comparación.
}\end{quote}


\section{Inteligencia Artificial en Educación}

El concepto de alfabetización en IA ha emergido como un marco para definir qué debe saber un ciudadano sobre inteligencia artificial. \citet{long2020ailiteracy} la definieron como el conjunto de competencias que permiten evaluar críticamente las tecnologías de IA, comunicarse con sistemas de IA y usarla como herramienta. \citet{ng2021ailiteracy} ampliaron la conceptualización e identificaron cuatro dimensiones: conocer la IA, usarla, evaluarla y abordar sus implicaciones éticas.

Existe una distinción relevante entre la IA como contenido de aprendizaje y la IA como herramienta pedagógica~\citep{holmes2019aied}. En el primer caso los estudiantes aprenden sobre qué es la IA y cómo funciona. En el segundo, los sistemas de IA apoyan procesos de enseñanza mediante tutores inteligentes o plataformas adaptativas. Una tercera dimensión, presente en las políticas más recientes, es la IA como objeto de regulación dentro del espacio educativo, donde la irrupción de los modelos generativos obligó a muchos sistemas a definir reglas de uso antes de poder incorporarla~\citep{unesco2023genai}.


\section{Marcos Internacionales}

Los organismos internacionales han configurado el discurso global sobre IA y educación. Sus marcos no son vinculantes, pero establecen vocabularios y prioridades que influyen en las estrategias nacionales. En esta investigación son relevantes como posibles fuentes de difusión.

La UNESCO fue pionera con el Consenso de Beijing de 2019, que estableció cinco áreas prioritarias~\citep{unesco2019beijing}. Después publicó una recomendación sobre ética de la IA~\citep{unesco2021ai} y una guía para IA generativa en educación~\citep{unesco2023genai}. \citet{miao2021guidance} tradujeron estos principios en recomendaciones concretas para sistemas educativos.

La OCDE, por su parte, adoptó en 2019 sus Principios sobre IA~\citep{oecd2019ai} y desarrolló el \textit{Education Policy Outlook}~\citep{oecd2021epo}, un marco estandarizado que informa las dimensiones del análisis de esta investigación.


\section{Políticas Públicas Educativas}

\citet{rizvi2010globalizing} definen las políticas educativas como el conjunto de directrices, regulaciones y asignaciones de recursos mediante las cuales los gobiernos organizan y orientan los sistemas de enseñanza. A diferencia de otros ámbitos de la política social, las políticas educativas forman capital humano, promueven cohesión social y transmiten cultura al mismo tiempo. Un rasgo recurrente de las políticas de IA en educación es su carácter programático más que legislativo: la mayoría de los documentos analizados son estrategias, planes de acción o lineamientos que expresan intenciones gubernamentales sin contar con mecanismos vinculantes.

Los países del corpus (es decir, el conjunto sistemático de documentos que constituyen la fuente de datos del análisis) se encuentran en etapas distintas del proceso de políticas. La literatura estándar en política pública~\citep{howlett2003studying} distingue fases sucesivas: establecimiento de agenda, formulación, adopción, implementación y evaluación. China, con siete documentos entre 2017 y 2025, se encuentra en fases avanzadas de implementación y evaluación. Sudáfrica, con un informe de comisión, aún se encuentra en establecimiento de agenda. Esta asimetría exige cautela al comparar documentos que capturan momentos distintos del proceso.


\section{Transferencia y Difusión de Políticas}

La globalización ha intensificado los flujos de ideas sobre política educativa entre países. \citet{dolowitz2000learning} propusieron un marco para entender la transferencia como un proceso en el que conocimientos de un sistema político se usan para desarrollar políticas en otro. En el campo educativo, esta transferencia rara vez es una copia literal. \citet{steinerkhamsi2014cross} distingue entre tres explicaciones para la similitud entre políticas: el préstamo genuino, la influencia común de un tercer actor y la convergencia funcional. A su vez, \citet{shipan2012diffusion} identificaron cinco mecanismos de difusión: aprendizaje, competencia, emulación, coerción y contagio social. Estos mecanismos no son excluyentes.

Cuando el análisis identifica alta similitud entre los documentos de dos países, estas explicaciones deben evaluarse. Dos países de familias lingüísticas diferentes que muestran alta similitud y ambos citan el Consenso de Beijing sugieren influencia común como hipótesis más parsimoniosa. En cambio, la alta similitud entre países angloparlantes obliga a considerar también el componente lingüístico. El diseño metodológico incluye un control lingüístico que permite descomponer estas fuentes.


\section{Educación Comparada}

La educación comparada es la disciplina que estudia los sistemas educativos de distintos países para entender sus similitudes, diferencias y las razones detrás de ambas. Su pertinencia para esta investigación es directa: comparar las políticas de IA en educación de siete países es, en esencia, un ejercicio de educación comparada. La disciplina ofrece los principios metodológicos que orientan qué se puede comparar, en qué nivel y con qué cautelas.

El marco tridimensional de \citet{bray1995levels} cruza niveles geográficos, grupos demográficos y aspectos de la educación. La presente investigación opera en el nivel nacional, abarca siete países y se enfoca en las políticas de IA en educación. Las nueve dimensiones de comparación que se presentan más adelante desagregan este aspecto en componentes específicos.

Una distinción útil para este trabajo es la que propusieron \citet{rizvi2010globalizing} entre política como texto (lo que el documento dice), como discurso (cómo lo enmarca) y como efecto (qué ocurre cuando se implementa). El análisis opera en los planos de texto y discurso. Las dimensiones de proceso y resultados capturan lo que los documentos dicen sobre procesos y metas, no lo que ocurre en la práctica.

Los métodos de análisis comparativo en este campo han sido predominantemente cualitativos, y esta investigación aporta herramientas computacionales en esa dirección.


\section{Herramientas de Procesamiento de Lenguaje Natural}

El análisis computacional de textos permite operar sobre un corpus multilingüe a una escala que la lectura humana individual no alcanza. \citet{grimmer2013text} establecieron el principio rector: los métodos computacionales amplifican la capacidad del investigador pero no reemplazan su juicio. La lectura humana sigue siendo indispensable para interpretar qué significan los patrones que las herramientas detectan. Pero las herramientas hacen posible algo que un investigador individual no puede hacer: comparar sistemáticamente fragmentos de texto en cuatro idiomas distintos y medir cuánto se parecen entre sí.

La técnica central es la representación semántica de textos: cada fragmento de un documento se convierte en una representación numérica que captura su significado. Un modelo de análisis multilingüe genera estas representaciones de manera que textos con significado similar quedan próximos, independientemente del idioma en que estén escritos~\citep{reimers2019sentencebert}. Esta capacidad multilingüe es la que permite comparar documentos originalmente escritos en inglés, español, alemán y chino dentro de un mismo marco de análisis. A pesar de esto, la separación entre idiomas no es perfecta~\citep{conneau2020unsupervised}, lo que motiva el control lingüístico.

El descubrimiento automático de temas identifica la estructura temática del corpus sin imponer categorías previas~\citep{grootendorst2022bertopic}. El método agrupa los fragmentos por similitud semántica y luego identifica las palabras más representativas de cada grupo. En esta investigación, el descubrimiento de temas se ejecuta antes de aplicar el marco de nueve dimensiones, siguiendo el principio de que primero se descubren patrones y luego se comparan con las categorías del investigador~\citep{nelson2020computational}. Los temas que no corresponden a ninguna dimensión son puntos ciegos del marco. Las dimensiones sin correspondencia son estructura impuesta. Ambas divergencias son hallazgos, no defectos del diseño.

Para aislar el efecto del idioma se implementan tres estrategias: una línea base monolingüe que traduce todos los documentos al inglés y recalcula la similitud, la comparación de similitudes entre países del mismo idioma versus países de idiomas diferentes, y un test con un documento disponible en múltiples idiomas.

\section{Marco Analítico}

El diseño de esta investigación requería un marco analítico que cumpliera tres condiciones: que fuera global (aplicable a políticas de países de distintas regiones), que fuera específico para IA en educación (no para IA en general), y que fuera operacionalizable computacionalmente (susceptible de convertirse en representaciones numéricas para el análisis semántico). No se encontró ningún marco existente que cumpliera las tres condiciones simultáneamente. El \textit{Education Policy Outlook} de la OCDE~\citep{oecd2021epo} ofrece categorías globales pero no es específico para IA. El Consenso de Beijing~\citep{unesco2019beijing} es específico para IA en educación pero no está diseñado para comparación operacionalizable. La solución fue derivar un marco propio a partir de estas fuentes, con una cadena explícita que va desde el marco teórico hasta la representación numérica.

Las nueve dimensiones resultantes se organizan en tres niveles que reflejan distintos planos de la política educativa.

El primer nivel agrupa las dimensiones de contenido, que describen qué temas abordan las políticas. La dimensión de \textbf{currículo} captura qué se enseña sobre IA: si es una materia independiente, una herramienta transversal o un tema de reflexión crítica; en qué niveles educativos se introduce y qué competencias define. La \textbf{formación docente} abarca los programas de capacitación inicial y continua, las competencias que se exigen a los profesores para trabajar con IA en el aula y los marcos pedagógicos que orientan esa formación. La \textbf{infraestructura} se refiere a la inversión en conectividad, equipamiento, plataformas digitales y servicios tecnológicos que hacen posible la implementación. La \textbf{investigación} cubre los centros de I+D, los programas de financiamiento, los vínculos entre universidades e industria y las prioridades de investigación declaradas.

El segundo nivel agrupa las dimensiones de proceso, que describen cómo se formulan las políticas. La \textbf{gobernanza} identifica qué ministerio u organismo lidera la política, cómo se coordinan los distintos niveles de gobierno y qué marcos regulatorios se establecen. La \textbf{participación de actores} examina los mecanismos de consulta pública, las alianzas público-privadas y la participación formal de sociedad civil, academia y sindicatos. Dos políticas pueden coincidir en los temas que abordan pero diferir radicalmente en quién las diseñó y quién tuvo voz en el proceso.

El tercer nivel agrupa las dimensiones orientadas a resultados, que capturan lo que las políticas prometen lograr y cómo lo verifican. La \textbf{ética y rendición de cuentas} incluye directrices sobre transparencia algorítmica, protección de datos, prevención de sesgos, supervisión humana y mecanismos de auditoría. La \textbf{equidad} abarca las medidas para reducir brechas digitales, atender a poblaciones vulnerables y prevenir que la IA reproduzca desigualdades. Las \textbf{metas y evaluación} comprenden los objetivos cuantificables, los indicadores de desempeño, los plazos de cumplimiento y los mecanismos de monitoreo.

La Tabla~\ref{tab:dimensiones} resume las nueve dimensiones con su nivel, fuente y lo que captura en el texto de las políticas.

\begin{table}[htbp]
\centering
\caption{Nueve dimensiones del marco analítico}
\label{tab:dimensiones}
\small
\begin{tabular}{llll}
\hline
\textbf{Dimensión} & \textbf{Nivel} & \textbf{Fuente} & \textbf{Captura en el texto} \\
\hline
Currículo & Contenido & OECD EPO; Beijing & Qué y cómo se enseña sobre IA \\
Formación docente & Contenido & OECD EPO; Beijing & Capacitación y competencias docentes \\
Infraestructura & Contenido & OECD EPO; Beijing & Conectividad, equipamiento, plataformas \\
Investigación & Contenido & OECD EPO; Beijing & I+D, vínculos universidad-industria \\
\hline
Gobernanza & Proceso & OECD EPO & Liderazgo institucional, coordinación \\
Participación & Proceso & Bray \& Thomas; Shipan & Consulta pública, alianzas, actores \\
\hline
Ética & Resultados & UNESCO Ethics; OECD AI & Transparencia, privacidad, auditoría \\
Equidad & Resultados & OECD EPO; Beijing & Brechas digitales, poblaciones vulnerables \\
Metas & Resultados & OECD EPO & Indicadores, plazos, revisión \\
\hline
\end{tabular}
\end{table}

\subsection{Operacionalización: de la dimensión al número}

Cada dimensión se operacionaliza mediante un párrafo ancla de 3 a 5 oraciones que describe qué se observa cuando un documento de política aborda esa dimensión. Estos párrafos no son extractos de los documentos del corpus: son textos escritos por el investigador a partir de las definiciones del marco teórico, antes de analizar los resultados. Un modelo de análisis semántico multilingüe convierte cada párrafo ancla en una representación numérica~\citep{reimers2019sentencebert}. La afinidad temática entre esta representación y cada segmento del corpus produce una puntuación por dimensión.

Los párrafos ancla se escriben en español e inglés, y la puntuación final promedia ambas versiones para reducir el sesgo que introduciría un solo idioma. Los párrafos se registran en el código del proyecto antes de ejecutar el análisis supervisado, como mecanismo de pre-registro~\citep{nosek2018preregistration}. Los procedimientos completos se detallan en el Capítulo~\ref{cap:metodologia}.
